%=============================================================================
% SECTION 1: INTRODUCTION
% DFD Unified Review --- Part I: Foundations
%=============================================================================

\section{Introduction}\label{sec:introduction}

%-----------------------------------------------------------------------------
% §1.1 THE LANDSCAPE OF GRAVITY THEORIES
%-----------------------------------------------------------------------------
\subsection{The Landscape of Gravity Theories}\label{sec:landscape}

Einstein's general relativity (\GR) has withstood a century of experimental 
scrutiny with remarkable success~\cite{Will2014LRR,Will2018}. Solar system tests, 
binary pulsar timing, and gravitational wave observations all confirm \GR's 
predictions to extraordinary precision. Yet the theory's success comes at a 
cost: explaining astrophysical and cosmological observations requires postulating 
that $95\%$ of the universe's energy content consists of dark matter and dark 
energy---components that have never been directly detected despite decades of 
experimental effort~\cite{Bertone2018,Peebles2003}.

Astrophysical anomalies relative to \GR with visible matter alone form a 
remarkably coherent pattern. Spiral galaxy rotation curves are flat rather than 
Keplerian~\cite{Rubin1980}; low surface-brightness galaxies follow tight scaling 
relations~\cite{McGaugh2016RAR}; galaxy clusters require additional mass beyond 
their baryonic content~\cite{Zwicky1933}; and large-scale structure and supernova 
data point to late-time accelerated expansion~\cite{Riess1998,Perlmutter1999}. 
The dominant response has been the \LCDM paradigm, which retains \GR but 
postulates cold dark matter and a cosmological constant.

An alternative approach modifies gravity itself. Modified Newtonian Dynamics 
(\MOND) introduced a characteristic acceleration scale 
$\azero \sim 10^{-10}\,\si{m/s^2}$ governing the transition between Newtonian 
and deep-field behavior in galaxies~\cite{Milgrom1983,FamaeyMcGaugh2012}. 
Remarkably, this single parameter successfully predicts rotation curves, the 
baryonic Tully-Fisher relation, and the radial acceleration relation across 
galaxies spanning five decades in mass~\cite{Lelli2017}.

A striking and poorly understood coincidence is that $\azero$ is numerically 
close to the cosmic acceleration scale $a_\Lambda \sim cH_0$ inferred from 
the expansion rate~\cite{Milgrom1983}. This suggests a possible deep connection 
between galactic dynamics and cosmology that \LCDM treats as accidental.

\begin{table}[h!]
\caption{Comparison of approaches to the gravitational puzzle.}
\label{tab:gravity_comparison}
\begin{ruledtabular}
\footnotesize
\begin{tabular}{lccc}
Theory & Key Feature & Status & DM/DE? \\
\hline
\GR + \LCDM & Curved spacetime & Standard & Both \\
\MOND & $\mu$-crossover & Empirical & Replaces DM \\
$f(R)$ & Modified action & Various & Modified \\
TeVeS & Tensor-vector-scalar & Falsified$^a$ & --- \\
Brans-Dicke & Scalar-tensor & Constrained & Modified \\
\DFD & Optical index & This work & MOND + LPI \\
\end{tabular}
\end{ruledtabular}
\footnotesize{$^a$GW170817 speed constraint~\cite{GW170817_speed}.}
\end{table}

Scalar-tensor theories have proliferated as alternatives to 
\GR~\cite{CliftonEtAl2012,Joyce2015}. Brans-Dicke theory~\cite{BransDicke1961} 
introduced a dynamical scalar coupled to curvature. Bekenstein's Tensor-Vector-Scalar 
theory (TeVeS)~\cite{Bekenstein2004} attempted to provide a relativistic 
completion of \MOND but was falsified by the near-simultaneous arrival of 
gravitational waves and light from GW170817~\cite{GW170817_speed}. 
The $f(R)$ family~\cite{Sotiriou2010} modifies the Einstein-Hilbert action 
directly. Each approach faces its own challenges: additional parameters, 
instabilities, or conflict with precision tests.

The theory presented in this review---Density Field Dynamics (\DFD)---takes a 
different path. Rather than modifying \GR's geometric structure, \DFD posits 
that spacetime is fundamentally flat but contains a scalar field establishing 
an optical refractive index. This approach has historical precedent: in 1911-12, 
before completing general relativity, Einstein himself explored gravity as a 
variable speed of light~\cite{Einstein1911,Einstein1912}. Gordon in 1923 showed 
that electromagnetic wave propagation in a medium can be described by an 
effective ``optical metric''~\cite{Gordon1923}. \DFD makes this optical 
perspective foundational rather than emergent.

Table~\ref{tab:gravity_comparison} summarizes how \DFD relates to other approaches. 
The key distinction is that \DFD reproduces \GR's predictions where tested 
(solar system, gravitational waves, binary pulsars) while making specific, 
falsifiable predictions where not yet tested (laboratory LPI tests, clock 
anomalies, matter-wave phases).

%-----------------------------------------------------------------------------
% §1.2 CORE IDEA: GRAVITY AS AN OPTICAL MEDIUM
%-----------------------------------------------------------------------------
\subsection{Core Idea: Gravity as an Optical Medium}\label{sec:core_idea}

The central insight of \DFD is that gravity can be understood as a refractive 
medium. Just as light bends when passing through glass because of a spatially 
varying refractive index, light and matter in a gravitational field respond 
to a cosmically varying index $n = e^\psi$. This is not merely an analogy---it 
is the complete dynamical content of the theory.

The formulation rests on two postulates that constitute the 
\emph{Minimal Optical Equivalence} principle:

\paragraph{Postulate P1 (Light).}
In a broadband nondispersive window, electromagnetic waves propagate according 
to the eikonal of an effective optical metric
\begin{equation}
d\tilde{s}^2 = -\frac{c^2\,dt^2}{n^2(\mathbf{x},t)} + d\mathbf{x}^2, 
\qquad n(\mathbf{x},t) = e^{\psifield(\mathbf{x},t)}.
\label{eq:optical_metric}
\end{equation}
This is the Gordon-Perlick optical geometry statement~\cite{Gordon1923,Perlick2000}, 
grounding ray optics in wave theory with a single scalar field $\psifield$ 
determining the local refractive index.

\paragraph{Postulate P2 (Matter).}
Test bodies move under the conservative potential
\begin{equation}
\Phipot \equiv -\frac{c^2}{2}\psifield, 
\qquad \avec = \frac{c^2}{2}\nabla\psifield = -\nabla\Phipot,
\label{eq:matter_response}
\end{equation}
which fixes the weak-field normalization to match \GR's classic optical tests 
(light deflection factor of two, Shapiro delay coefficient, gravitational redshift).

The exponential form $n = e^\psi$ is not arbitrary but follows from three 
requirements:
\begin{enumerate}[label=(\roman*)]
\item \emph{Positivity:} $n > 0$ everywhere, ensuring light propagation is 
always defined.
\item \emph{Weak-field limit:} For $|\psi| \ll 1$, we have 
$n \approx 1 + \psi$, recovering the linear regime.
\item \emph{Multiplicative composition:} Sequential media combine as 
$n_{\text{total}} = n_1 n_2 = e^{\psi_1 + \psi_2}$, matching the additive 
nature of gravitational potentials.
\end{enumerate}

The factor-of-two deflection that matches \GR emerges automatically. 
In \GR, light deflection receives equal contributions from spatial curvature 
and time dilation. In \DFD, the optical metric~\eqref{eq:optical_metric} 
encodes both effects: the phase velocity $c/n$ slows in the potential well, 
and wavefronts tilt toward the slower region. The result is precisely 
$2GM/(c^2 b)$ at impact parameter $b$---the same as \GR.

% FIGURE 1: Conceptual comparison
\begin{figure}[t]
\centering
\begin{tikzpicture}[scale=0.9]
% Panel A: GR
\begin{scope}[xshift=-1.5cm]
\node at (0,2.5) {\textbf{(a) General Relativity}};
% Curved grid
\draw[gray,thin] (-1.8,-1.5) to[out=80,in=260] (-1.5,1.5);
\draw[gray,thin] (-0.6,-1.5) to[out=75,in=265] (-0.3,1.5);
\draw[gray,thin] (0.6,-1.5) to[out=105,in=255] (0.3,1.5);
\draw[gray,thin] (1.8,-1.5) to[out=100,in=260] (1.5,1.5);
\draw[gray,thin] (-2,-1) to[out=10,in=170] (2,-1);
\draw[gray,thin] (-2,0) to[out=5,in=175] (2,0);
\draw[gray,thin] (-2,1) to[out=-5,in=185] (2,1);
% Mass
\shade[ball color=orange!80!black] (0,0) circle (0.3);
% Geodesic
\draw[blue,thick,-{Stealth}] (-2,1.2) to[out=-10,in=170] (2,0.8);
\node[blue] at (0,-2) {Geodesic on curved manifold};
\end{scope}
% Panel B: DFD
\begin{scope}[xshift=3.5cm]
\node at (0,2.5) {\textbf{(b) Density Field Dynamics}};
% Flat grid
\draw[gray,thin] (-1.8,-1.5) -- (-1.8,1.5);
\draw[gray,thin] (-0.6,-1.5) -- (-0.6,1.5);
\draw[gray,thin] (0.6,-1.5) -- (0.6,1.5);
\draw[gray,thin] (1.8,-1.5) -- (1.8,1.5);
\draw[gray,thin] (-2,-1) -- (2,-1);
\draw[gray,thin] (-2,0) -- (2,0);
\draw[gray,thin] (-2,1) -- (2,1);
% Index gradient (shading)
\fill[orange!20,opacity=0.5] (0,0) circle (1.5);
\fill[orange!40,opacity=0.5] (0,0) circle (1.0);
\fill[orange!60,opacity=0.5] (0,0) circle (0.5);
% Mass
\shade[ball color=orange!80!black] (0,0) circle (0.3);
% Bent ray
\draw[red,thick,-{Stealth}] (-2,1.2) to[out=-5,in=175] (2,0.8);
\node[red] at (0,-2) {Ray bent by $n(\mathbf{x})$};
\end{scope}
\end{tikzpicture}
\caption{Conceptual comparison of (a)~General Relativity, where gravity curves 
spacetime and particles follow geodesics on a curved manifold, and 
(b)~Density Field Dynamics, where spacetime is flat but contains a refractive 
medium with index $n(\mathbf{x}) = e^{\psi(\mathbf{x})}$ that bends light rays.
Both yield identical weak-field predictions.}
\label{fig:concept}
\end{figure}

Figure~\ref{fig:concept} illustrates the conceptual difference. In \GR, 
gravity is geometry: mass curves spacetime, and particles follow geodesics 
on a curved manifold. In \DFD, spacetime remains flat (Minkowski background), 
but a scalar field creates a refractive medium. The observational predictions 
are identical in the weak-field regime---the theories differ only in their 
ontology and in specific strong-field or laboratory contexts.

The connection between the two postulates is not coincidental. Both light 
and matter respond to the same field $\psi$, ensuring the Weak Equivalence 
Principle is satisfied: all test masses fall with the same acceleration 
$\avec = (c^2/2)\nabla\psi$ regardless of composition. The universality of 
free fall is built into the structure.

%-----------------------------------------------------------------------------
% §1.3 WHAT DFD CLAIMS AND WHAT IT DOESN'T
%-----------------------------------------------------------------------------
\subsection{What DFD Claims and What It Doesn't}\label{sec:claims}

Before proceeding to the technical development, we state explicitly what 
\DFD claims and what it does not claim. This serves to preempt misinterpretation 
and to define the scope of falsifiability.

\paragraph{Claim taxonomy.}
For clarity, this review uses three claim types.  \textbf{Core-derived} statements follow from the DFD field equations and actions presented in the main text.  \textbf{Auxiliary-closure-derived} statements follow from explicitly displayed supplemental structural postulates (for example, the finite-symmetry closure used in the microsector).  \textbf{Empirical consistency} statements summarize benchmark calculations, fits, or data confrontations.  This taxonomy is used to keep the one-paper presentation logically unified without blurring the difference between core theorems, closure-framework consequences, and benchmark evidence.

\paragraph{What \DFD Claims:}
\begin{enumerate}[label=\arabic*.]
\item \textbf{Weak-field equivalence with \GR:} The optical metric with 
$n = e^\psi$ reproduces all Solar System tests. The Parametrized Post-Newtonian 
(PPN) parameters are $\gamma = \beta = 1$, and all ten PPN parameters match 
\GR at first post-Newtonian order (\S\ref{sec:ppn}).

\item \textbf{Gravitational waves at speed $c$:} A minimal transverse-traceless 
sector propagates at the speed of light with two tensor polarizations, 
consistent with GW170817 and LIGO/Virgo/KAGRA observations (\S\ref{sec:gw}).

\item \textbf{\MOND-like phenomenology:} At galactic scales where 
$|\nabla\psi|/\astar \ll 1$, a nonlinear crossover function $\mu(x)$ produces 
flat rotation curves, the baryonic Tully-Fisher relation, and the radial 
acceleration relation without cold dark matter (\S\ref{sec:galactic}).

\item \textbf{Channel-resolved clock and cavity residuals:} \DFD predicts that clock responses are channel-dependent rather than universal. Same-ion optical ratios tightly constrain the pure $\alpha$ sector, cross-species and nuclear clocks probe composition and strong-sector channels, and cavity--atom comparisons reduce at tree level to a screened residual rather than an order-unity slope (\S\ref{sec:clocks}, \S\ref{sec:cavity}).

\item \textbf{Matter-wave $T^3$ signature:} Atom interferometers should 
exhibit a small $T^3$ contribution to the phase proportional to $\nabla|\nabla\psi|$, 
absent in \GR at leading order (\S\ref{sec:matter_wave}).

\item \textbf{Parameter-free $\alpha$-relations:} Three numerical coincidences 
link the fine-structure constant $\alpha$ to gravitational scales without 
free parameters:
\begin{align}
\azero &= 2\sqrt{\alpha}\,cH_0, \\
\ka &= 3/(8\alpha) \approx 51.4, \\
\kalpha &= \alpha^2/(2\pi) \approx 8.5 \times 10^{-6}.
\end{align}
The first predicts the \MOND acceleration scale to within 3\%; the second 
and third enter clock phenomenology (\S\ref{sec:alpha}).

\item \textbf{CMB from pure $\psi$-physics:} The CMB peak structure is derived 
directly from $\psi$-physics without dark matter. Peak ratio $R \approx 2.4$ 
arises from baryon loading in $\psi$-gravity; peak location $\ell_1 \approx 220$ 
arises from $\psi$-lensing (gradient-index optics with $n = e^\psi$). 
\textbf{Quantitative reconstruction:} $\Delta\psi(z=1) = 0.27 \pm 0.02$ from $H_0$-independent 
distance ratios explains the ``accelerating expansion'' as an optical effect.
\textbf{No dark matter; no dark energy; one cosmological screen $\Delta\psi$} (\S\ref{sec:psi-reconstruction}).
\end{enumerate}

\paragraph{Theoretical Completeness :}
\begin{enumerate}[label=\arabic*.]
\item \textbf{UV completion from topology:} The $\mathbb{C}P^2 \times S^3$ gauge emergence framework provides UV completion. Unlike GR, DFD has flat spacetime (no curvature singularities) and classical $\psi$ (action $\gg \hbar$). The topology derives all ``constants''---this IS the UV physics (\S\ref{sec:open-uv}).

\item \textbf{CMB derived analytically:} Peak ratio $R = 2.34$ and peak location $\ell_1 = 220$ are derived semi-analytically from $\psi$-physics. CLASS/CAMB are GR-based tools; the DFD derivation is complete without them.

\item \textbf{Cluster mechanism RESOLVED:} Multi-scale averaging + updated baryonics yields Obs/DFD $= 0.98 \pm 0.05$ for all 16 clusters (100\% within $\pm$10\%). Galaxy groups show EFE suppression as predicted (\S\ref{sec:cosmo-clusters}, Appendix~\ref{app:cluster-full}).

\item \textbf{Standard Model from topology:} The gauge emergence framework 
(\S\ref{sec:quantum}) derives: $SU(3) \times SU(2) \times U(1)$ from $(3,2,1)$ partition,
$N_{\rm gen} = 3$ from index theory, $\alpha = 1/137$ from Chern-Simons, 
all 9 charged fermion masses (1.42\% error), CKM and PMNS matrices, 
$v = M_P \alpha^8 \sqrt{2\pi}$ (hierarchy solved), and $\bar\theta = 0$ to all orders (Theorem~\ref{thm:strong_cp_all_loops}; no axion required).
Physical validity conditional on DFD gravity being confirmed experimentally.

\item \textbf{Scope boundary:} Loop corrections in the $\psi$-gauge coupled system are not computed; the classical/EFT level is sufficient for all predictions.
\end{enumerate}

The philosophy is: \emph{conservative where tested, bold where testable}. 
\DFD reproduces \GR in all regimes where \GR has been confirmed, and makes 
specific, quantitative predictions in regimes where decisive tests are 
experimentally accessible.

%-----------------------------------------------------------------------------
% §1.4 READER'S GUIDE
%-----------------------------------------------------------------------------
\subsection{Reader's Guide}\label{sec:reader_guide}

This review is organized to be readable both linearly and as a reference. 
The structure follows a logical progression from foundations to frontiers, 
with each part addressing a distinct aspect of the theory.

\paragraph{Part I: Foundations (Sections~\ref{sec:introduction}--\ref{sec:wellposed}).}
Establishes the mathematical framework: the optical metric, action principle, 
field equations, and proof of well-posedness (existence, uniqueness, stability). 
This part is prerequisite for all subsequent sections.

\paragraph{Part II: Contact with Known Physics (Sections~\ref{sec:ppn}--\ref{sec:gw}).}
Demonstrates that \DFD reproduces \GR where tested. Section~\ref{sec:ppn} 
presents the complete PPN analysis showing $\gamma = \beta = 1$. 
Section~\ref{sec:gw} develops the gravitational wave sector and verifies 
consistency with LIGO/Virgo/KAGRA constraints.

\paragraph{Part III: Strong Fields (Section~\ref{sec:strong}).}
Extends to strong-field regimes: spherically symmetric solutions, photon 
spheres, optical horizons, and black hole shadows. Comparison with EHT 
observations of M87* and Sgr~A* is presented.

\paragraph{Part IV: Galactic Dynamics (Section~\ref{sec:galactic}).}
Develops the deep-field regime where $\mu \neq 1$: rotation curves, 
Tully-Fisher relation, and the radial acceleration relation. The single 
calibration on RAR data is described.

\paragraph{Part V: The $\alpha$-Relations (Section~\ref{sec:alpha}).}
Presents the three parameter-free numerical relations linking $\alpha$ to 
gravitational phenomenology, with derivation and verification.

\paragraph{Part VI: Laboratory Tests (Sections~\ref{sec:clocks}--\ref{sec:matter_wave}).}
Details the decisive experimental discriminators: atomic clock anomalies 
(\S\ref{sec:clocks}), cavity-atom LPI tests (\S\ref{sec:cavity}), 
and matter-wave interferometry (\S\ref{sec:matter_wave}). These sections 
are self-contained and can be read independently after Part~I.

\paragraph{Part VII: Frontiers and Open Problems 
(Sections~\ref{sec:cosmology}--\ref{sec:conclusions}).}
Addresses cosmological implications (\S\ref{sec:cosmology}), the conditional 
quantum/gauge sector (\S\ref{sec:quantum}), open problems and limitations 
(\S\ref{sec:openproblems}), and conclusions (\S\ref{sec:conclusions}).

\paragraph{Dependencies.}
\begin{itemize}
\item Sections~\ref{sec:introduction}--\ref{sec:wellposed} (Part~I) are 
prerequisite for all subsequent sections.
\item Section~\ref{sec:ppn} (PPN) is independent of galactic phenomenology 
(Section~\ref{sec:galactic}).
\item Laboratory tests (Sections~\ref{sec:clocks}--\ref{sec:matter_wave}) 
require only Part~I.
\item Strong fields (Section~\ref{sec:strong}) requires 
Sections~\ref{sec:formalism}--\ref{sec:wellposed}.
\end{itemize}

\paragraph{Notation.}
Standard notation is defined in Appendix~\ref{app:notation} and summarized 
here. The scalar field is $\psi$; the refractive index is $n = e^\psi$; 
the acceleration is $\avec = (c^2/2)\nabla\psi$; the crossover function 
is $\mu(x)$ with $x = |\nabla\psi|/\astar$; the acceleration scale is 
$\azero \sim 10^{-10}\,\si{m/s^2}$. Key equations are numbered sequentially 
throughout; a summary table appears in Appendix~\ref{app:derivations}.

\paragraph{A note on falsifiability.}
Every scientific theory must specify conditions under which it would be 
falsified. For \DFD, the decisive tests are:
\begin{itemize}
\item \textbf{Channel-resolved clocks and cavity residuals:} If same-ion, cross-species, and nuclear-clock data cannot be organized by the channel-resolved structure of Eq.~\eqref{eq:K-general}, the present clock mechanism is wrong. In particular, a high-precision null in cross-species atomic ratios and in the surviving $^{229}$Th/Sr nuclear window would remove the leading live laboratory channels.
\item \textbf{Cavity--atom residuals:} After geometric cancellation, the cavity--atom observable is no longer an order-unity discriminator but a screened residual. A future dedicated null at the residual sensitivity target of Sec.~\ref{sec:cavity} would constrain or remove that channel; a null only at the old $\delta\xi_{\rm LPI}<0.1$ level would not.
\item \textbf{Gravitational waves:} If ppE parameters deviate from zero 
in the strong-field regime, the radiative sector requires modification.
\end{itemize}
The theory is constructed to be falsifiable, not merely 
``not yet falsified.''

%-----------------------------------------------------------------------------
% §1.5 ASSUMPTIONS AND DEGREES OF FREEDOM LEDGER
%-----------------------------------------------------------------------------
\subsection{Assumptions and Degrees of Freedom Ledger}\label{sec:dof_ledger}

To prevent any accusation of hidden parameter tuning, we provide an explicit accounting of all inputs, outputs, and falsifiers. This ``ledger'' makes the theory's structure transparent.

\begin{table}[h]
\centering
\caption{Complete accounting of DFD inputs, outputs, and falsifiers.}\label{tab:dof_ledger}
\begin{ruledtabular}
\scriptsize
\begin{tabular}{lll}
\textbf{Category} & \textbf{Item} & \textbf{Status} \\
\hline
\multicolumn{3}{l}{\textit{Foundational Postulates (2)}} \\
& $n = e^\psi$ & Postulate \\
& $\Phi = -c^2\psi/2$ & Postulate \\
\hline
\multicolumn{3}{l}{\textit{Topological Data (from SM)}} \\
& $q_1 = 3$ & From SM \\
& $n = 5$ (multiplets) & SM def. \\
& $(a,n) = (9,5)$ & Unique \\
& $k_{\max} = 60$ & Bundle \\
& $N_{\rm gen} = 3$ & Index thm. \\
\hline
\multicolumn{3}{l}{\textit{Scale Input (1 measurement)}} \\
& $H_0$ or $G$ & Measured \\
\hline
\multicolumn{3}{l}{\textit{Functional Choice}} \\
& $\mu(x)$ form & Discrete \\
\hline
\multicolumn{3}{l}{\textit{Derived (0 free parameters)}} \\
& $\alpha^{-1} = 137.036$ & CS quant. \\
& $a_0 = 2\sqrt{\alpha}cH_0$ & Derived \\
& $G\hbar H_0^2/c^5 = \alpha^{57}$ & Derived \\
& $v = M_P\alpha^8\sqrt{2\pi}$ & Derived \\
& Masses, CKM, PMNS & Derived \\
\hline
\multicolumn{3}{l}{\textit{Falsifiers}} \\
& Cavity--atom residual null & Cavity \\
& Clock channel structure fails & Clocks \\
& $c_T \neq c$ & GW \\
& RAR $>3\sigma$ off & Galactic \\
\end{tabular}
\end{ruledtabular}
\end{table}

\paragraph{Key point.}
The $\mu(x)$ crossover function is \textbf{not} a continuous fit parameter. Its single scale $a_0$ is derived from the $\alpha$-relation $a_0 = 2\sqrt{\alpha}\,cH_0$; the functional form $\mu(x) = x/(1+x)$ is \textbf{uniquely determined} by the $S^3$ Chern-Simons microsector topology (Appendix~\ref{app:mu-derivation}). Once $H_0$ is measured, no adjustable parameters remain.

\paragraph{Clarification: Parameter structure.}
DFD has: (i)~zero continuous fit parameters analogous to $\Omega_m$, $w$, or CDM concentrations; (ii)~two topological integers ($k_{\max} = 60$, $N_{\rm gen} = 3$); (iii)~one empirical scale ($H_0$ or equivalently $G$). The Planck vs SH0ES tension in $H_0$ ($67.4 \pm 0.5$ vs $73.0 \pm 1.0$ km/s/Mpc) propagates to a corresponding $\sim$8\% range in $a_0$ predictions. Given any specific $H_0$ value, all $\alpha$-relations become predictions, not fits.
